% ruby-basic.tex -- P10: ruby.sty provider adapter, yoko placement + boundary glue.
% 自己表明型 (RUBY-PASS/RUBY-FAIL をログへ出力; runner の T9 ゲートが grep)。
% 検査対象（plan Approach 10 + Verification の ruby 行）:
%   1. 横組み: ルビは本文の上 (annotation above base)
%   2. 境界グルー: 、→ルビ の 0.5zw ペアグルーが 1 回だけ出る
%      (本文直接組との差分 0 = 二重発行なし・欠落なし)
%   3. ルビ→次の和文 の境界 (base 右端クラスで評価)
%   4. ルビ→欧文 の右エッジ xkanjiskip (0.25zw)
%   5. 欧文→ルビ の左エッジ xkanjiskip
%   6. 連続ルビ: overlap の前ルビふち出し kern (\ruby@width) が生きる
%   7. 本文より長いルビ: nooverlap で外枠がルビ幅になる
% 測定方法: ペア位置は \pdfsavepos ではなく box 幅差分
%   (wd(ルビを含む行片) - wd(対応する直接組)) == 0 で境界グルー同一。
\documentclass[pdflatex,paper=a5]{jlreq}
\pagestyle{empty}
\usepackage{ruby}
\makeatletter
\newbox\rboxa
\newbox\rboxb
\newdimen\rdiff
% RTEST <name> <tolerance> {ruby-side} {plain-side}:
%   2 つの hbox 幅の差が tol 以内なら RUBY-PASS。
\long\def\RTEST#1#2#3#4{%
  \setbox\rboxa=\hbox{#3}% 前後の文脈ごと凍結
  \setbox\rboxb=\hbox{#4}%
  \rdiff=\dimexpr\wd\rboxa-\wd\rboxb\relax
  \ifdim\rdiff<\z@ \rdiff=-\rdiff \fi
  \ifdim\rdiff<#2\relax
    \typeout{RUBY-PASS: #1 (diff=\the\rdiff)}%
  \else
    \typeout{RUBY-FAIL: #1 (diff=\the\rdiff, ruby=\the\wd\rboxa, plain=\the\wd\rboxb)}%
  \fi}
% 幅だけの比較では「ルビ箱が外にはみ出す」場合を区別できないため、
% 基準側にも同じ幅の phantom を置く (\rlap で重ねた実組でも可)。
% --- ここでの比較は「ルビ両側の境界グルー」に限る: ルビ本体幅は
% provider 計算 (base/annotation 幅) に従うため、基準側に
% \ruby と同じ box を置いて両側だけを交換する。
\begin{document}
% ================= 1. 横: annotation は上 =================
% annotation の baseline 位置 = base の上 (provider の vbox:
%   .7\baselineskip - .5ex + .3 ann-baselineskip)。実測で
%   annotation の深さ方向位置が base baseline より上であることを
%   確認する: ルビ全体 box の高さ > base だけの box の高さ。
\setbox\rboxa=\hbox{\ruby{漢}{かん}}%
\setbox\rboxb=\hbox{漢}%
\ifdim\ht\rboxa>\ht\rboxb
  \typeout{RUBY-PASS: annotation-above (ht=\the\ht\rboxa\space vs base \the\ht\rboxb)}%
\else
  \typeout{RUBY-FAIL: annotation-above (ht=\the\ht\rboxa\space vs base \the\ht\rboxb)}%
\fi
% annotation の横位置: base グループの中央 (hss センタリング)。
%   ann box が base 幅より狭い → ann の左右空きは対称。
\setbox\rboxa=\hbox{\ruby{漢}{か}}%
% \box0 (provider の base box) は \ruby 呼び出し後も残る:
%   ann box は \@tempboxa (消費される)。代わりに幾何で確認:
%   ルビ box 幅 = base 幅 (wd0; ann が狭い場合)。
\ifdim\wd\rboxa=\dimexpr1\jlreq@jareg\relax
  \typeout{RUBY-PASS: narrow-ann-box-width (wd=\the\wd\rboxa)}%
\else
  \typeout{RUBY-FAIL: narrow-ann-box-width (wd=\the\wd\rboxa, expected \the\jlreq@jareg)}%
\fi
% ================= 2. 境界グルー: 、→ルビ =================
% \ruby{字}{じ} の base を plain の 字 に置き換えた対照行と
% 全体幅が一致すれば、境界グルーが同一 (かつ 1 回だけ)。
\RTEST{punct-left-glue}{0.02pt}{漢、\ruby{字}{じ}あ}{漢、字あ}
% ルビの右: base 右端→次の和文
\RTEST{right-glue-cjk}{0.02pt}{漢\ruby{字}{じ}あ}{漢字あ}
% ルビの右→欧文 (xkanjiskip 0.25zw)
\RTEST{right-glue-latin}{0.02pt}{漢\ruby{字}{じ}a}{漢字a}
% 欧文→ルビの左 (xkanjiskip 0.25zw)
\RTEST{left-glue-latin}{0.02pt}{a\ruby{字}{じ}b}{a字b}
% 区切り約物の左エッジ (「、」はルビ base に入る)
\RTEST{punct-base-inside}{0.02pt}{漢\ruby{字、}{じ}あ}{漢字、あ}
% ================= 3. 連続ルビ (overlap) =================
% 2 つ目のルビは 1 つ目の ruby@width (= (ann幅-base幅)/2 = 0.9zw) を
% kern(\@tempdimb) として進む (hmode では先頭ふち出し kern は出ない:
% provider の vm kern は真の段落頭のみ)。全体幅 = 10 + 9 + 10。
\setbox\rboxa=\hbox{\ruby{漢}{かんたんいそん}\ruby{字}{じ}}%
\rdiff=\wd\rboxa
\advance\rdiff-\dimexpr2\jlreq@jareg+0.9\jlreq@jareg\relax
\ifdim\rdiff<\z@ \rdiff=-\rdiff \fi
\ifdim\rdiff<0.01pt
  \typeout{RUBY-PASS: consecutive-overlap-advance (wd=\the\wd\rboxa)}%
\else
  \typeout{RUBY-FAIL: consecutive-overlap-advance (wd=\the\wd\rboxa, expected 2.9\jlreq@jareg, residual=\the\rdiff)}%
\fi
% ================= 4. 本文より長いルビ (nooverlap) =================
\rubynooverlap
\setbox\rboxa=\hbox{\ruby{漢}{かんたんいそんですよ}}%
% 外枠 = ann 幅 (10 glyphs × 0.4 × f@size; f@size=9.99994pt なので
% 厳密には 39.99939pt — 0.01pt の許容で検査)。
\rdiff=\wd\rboxa
\rdiff=-\rdiff
\advance\rdiff4\jlreq@jareg
\ifdim\rdiff<\z@ \rdiff=-\rdiff \fi
\ifdim\rdiff<0.01pt
  \typeout{RUBY-PASS: long-ann-nooverlap-width (wd=\the\wd\rboxa)}%
\else
  \typeout{RUBY-FAIL: long-ann-nooverlap-width (wd=\the\wd\rboxa, expected 4\jlreq@jareg, diff=\the\rdiff)}%
\fi
% base が長い場合: 外枠 = base 幅。
\setbox\rboxa=\hbox{\ruby{依存関係}{いぞん}}%
\ifdim\wd\rboxa=\dimexpr4\jlreq@jareg\relax
  \typeout{RUBY-PASS: long-base-nooverlap-width (wd=\the\wd\rboxa)}%
\else
  \typeout{RUBY-FAIL: long-base-nooverlap-width (wd=\the\wd\rboxa, expected \the\dimexpr4\jlreq@jareg\relax)}%
\fi
\rubyoverlap
% ================= 5. provider パラメータ =================
% rubysize: annotation のフォントサイズ = 0.5\f@size。
\renewcommand{\rubysize}{0.5}%
\setbox\rboxa=\hbox{\ruby{漢}{かん}}%
% ann 幅 = 2 × 0.5 × 10pt = 10pt (= base 幅 → ruby@width = 0)
\ifdim\wd\rboxa=\jlreq@jareg
  \typeout{RUBY-PASS: rubysize-05 (wd=\the\wd\rboxa)}%
\else
  \typeout{RUBY-FAIL: rubysize-05 (wd=\the\wd\rboxa, expected \the\jlreq@jareg)}%
\fi
\renewcommand{\rubysize}{0.4}%
% rubysep: \rubysep は provider が kern として使うだけなので
% box の高さに反映される (テスト: 0 にすると kern が消える)。
\renewcommand{\rubysep}{0pt}%
\setbox\rboxa=\hbox{\ruby{漢}{かん}}%
\setbox\rboxb=\hbox{\ruby{漢}{かん}}%
\ifdim\ht\rboxa=\ht\rboxb
  \typeout{RUBY-PASS: rubysep-set (ht=\the\ht\rboxa)}%
\else
  \typeout{RUBY-FAIL: rubysep-set (\the\ht\rboxa\space vs \the\ht\rboxb)}%
\fi
\renewcommand{\rubysep}{-0.5ex}%
% ================= 6. 段落先頭・末尾 =================
% 段落先頭ルビ: 行頭グルーなし (parbdd 90→first は 0-natural)。
% コンパイル成立 + 行頭位置で検査 (PDF 側は runner の幾何ゲート)。
段落先頭のルビ。
\par
段落末尾の\ruby{ルビ}{ruby}
\par
% ================= 7. 書体・サイズ変更 =================
{\gtfamily\ruby{漢}{かん}ゴシック}
{\fontsize{14pt}{14pt}\selectfont\ruby{漢}{かん}大きい}
\ruby{\gtfamily 漢}{かん}
\typeout{RUBY-PASS: font-change-compiled}%
% ================= 8. 状態隔離: ルビの後の境界 =================
% ルビ直後の次の和文境界が base 右端クラスで評価されること。
% (ルビの base が句読点で終わる場合: pair 7→0 = 0.5zw が出る)
\RTEST{state-exit-baseclass}{0.02pt}{\ruby{漢、}{かん}字}{漢、字}
% 逆: base が開き括弧で始まる (pair 90→1 = 0)
\RTEST{state-line-head-open}{0.02pt}{\ruby{「あ」}{かん}い}{「あ」い}
% ルビ内で空きなく続く base の内部境界 (箱の中は隔離フレーム)
\RTEST{base-internal-isolated}{0.02pt}{x\ruby{漢、}{かん}x}{x漢、x}
% ================= 9. rubylatin (provider の latin モード) ====
% \rubylatin は境界グルーを出さない契約: 基準 = グルーなし直接組。
\rubylatin
\RTEST{rubylatin-no-edge-glue}{0.02pt}{漢\ruby{字}{じ}あ}{漢字あ}
\rubyCJK
\typeout{RUBY-DONE}
\end{document}
